\section{Заключение}\label{sec:conclusion}

\subsection{Достигнутые результаты}

В работе исследована структура границы обучаемости однослойной нейронной сети
в многомерном пространстве гиперпараметров, включающем шаг обучения, инерцию,
затухание весов, размер мини-батча и параметры адаптивных оптимизаторов.
Получены следующие результаты.

\begin{enumerate}[leftmargin=*,topsep=2pt,itemsep=3pt]
  \item \textbf{Теоретический мини-результат.} Для квадратичной функции
        $L(w)=\tfrac12\lambda w^2$ градиентный спуск с инерцией записан в виде
        двумерного линейного отображения с матрицей перехода
        $M = \bigl(\begin{smallmatrix} 1-\eta\lambda & -\eta\beta \\ \lambda & \beta \end{smallmatrix}\bigr)$,
        для которой $\det M = \beta$ и $\operatorname{tr} M = 1+\beta-\eta\lambda$.
        Доказано (предложение~\ref{prop:stab}), что область устойчивости есть
        $0 < \eta\lambda < 2(1+\beta)$ при $|\beta|<1$, то есть инерция
        расширяет допустимый шаг обучения вплоть до двукратного по сравнению
        с обычным градиентным спуском; при комплексных собственных значениях
        скорость сходимости равна $\sqrt{\beta}$ и не зависит от кривизны.
        Результат подтверждён численно.
  \item \textbf{Репликация.} Воспроизведён базовый эксперимент предыдущей НИР;
        для активации tanh на сетке $512\times512$ получено
        $D_{\mathrm{box}} = \DBaseTanh \pm \DBaseTanhS$. Показано, что
        сравнение активаций требует индивидуального выбора окна сканирования,
        поскольку порог устойчивости у них различается на порядки.
  \item \textbf{Фрактальность сохраняется при добавлении инерции.} В сечении
        $(\eta,\beta)$ получено $D_{\mathrm{box}} = \DEtaBeta \pm \DEtaBetaS$,
        что совпадает с базовой картой в пределах погрешности.
  \item \textbf{Влияние регуляризации и батча.} Затухание весов сглаживает
        границу ($\DEtaWd$ против $\DSgd$ при том же разрешении), уменьшение
        размера батча --- увеличивает изрезанность монотонно:
        $\DMbFull$, $\DMbFour$, $\DMbOne$ при $B=32,4,1$.
  \item \textbf{Сравнение оптимизаторов.} Получены размерности $\DSgd$ (SGD),
        $\DMomentum$ (SGD с инерцией $\beta=0.9$), $\DAdam$ (Adam),
        $\DRmsprop$ (RMSprop). Гипотеза о сглаживающем действии адаптивных
        методов подтверждается лишь частично: сглаживание обеспечивает
        усреднение градиента по времени, а не покоординатная нормализация, и
        является свойством выбранного сечения --- в сечениях $(\eta,\beta_1)$
        и $(\eta,\beta_2)$ сам Adam даёт $\DAdamBOne$ и $\DAdamBTwo$.
  \item \textbf{Структурированные данные.} Переход от случайных данных к набору
        digits увеличивает размерность границы ($\DSynth \to \DDigits$);
        фрактальность не является артефактом синтетических данных.
  \item \textbf{Трёхмерное пространство.} Построены четыре сечения
        $(\eta,\beta)$ при разных $\lambda_{wd}$; зависимость размерности от
        $\lambda_{wd}$ немонотонна, что показывает ограниченность вывода о
        трёхмерной границе по набору двумерных сечений.
  \item \textbf{Воспроизводимый инструментарий.} Реализована ансамблевая
        векторизованная схема сканирования, позволяющая обсчитывать сетки
        $512\times512$ на одном ядре CPU за единицы минут, с фиксацией seed,
        точности и полным сохранением параметров каждого запуска.
\end{enumerate}

\subsection{Научная и практическая значимость}

Основной содержательный вывод состоит в том, что фрактальность границы
обучаемости --- свойство не конкретного оптимизатора, а самой нелинейной
динамики обучения вблизи порога устойчивости. Переход к более сложным
оптимизаторам меняет расположение границы и то, в каких координатах
проявляется её нерегулярность, но не устраняет её. Это уточняет ожидание,
согласно которому адаптивные методы «сглаживают» ландшафт обучаемости.

Практически это означает, что чувствительность к шагу обучения не может быть
устранена выбором оптимизатора и должна учитываться в процедурах
автоматического подбора гиперпараметров; при этом умеренная $L_2$-регуляризация
и увеличение размера батча объективно повышают предсказуемость обучения.

\subsection{Ограничения}

Основные ограничения перечислены в разделе~\ref{sec:limitations}. Кратко:
оценка размерности систематически зависит от разрешения сетки, что допускает
сравнение лишь внутри одного разрешения; окна сканирования подбирались
индивидуально, что делает межоптимизаторное сравнение качественным; переход
между режимами представляет собой протяжённую зону перемешивания, а не
кривую; критерий сходимости, унаследованный от предыдущей работы, засчитывает
как успех вырожденный случай «мёртвых» ReLU-нейронов; исследована лишь
однослойная сеть при $T=400$ итерациях.

\subsection{Направления дальнейшей работы}

\begin{enumerate}[leftmargin=*,topsep=2pt,itemsep=2pt]
  \item Прямой трёхмерный box-counting в пространстве
        $(\eta,\beta,\lambda_{wd})$ вместо набора сечений; это требует
        примерно двух порядков дополнительного машинного времени либо
        адаптивного разрешения сетки вблизи границы.
  \item Уточнение критерия обучаемости: замена порога «loss ниже начального»
        на требование снижения loss в заданное число раз с отдельной
        классификацией вырожденных состояний (мёртвые нейроны, стягивание
        весов к нулю регуляризацией).
  \item Проверка устойчивости выводов к числу итераций: измерение
        $D_{\mathrm{box}}(T)$ и определение, сходится ли оценка при $T\to\infty$.
  \item Перенос сравнения оптимизаторов на глубокие архитектуры и на
        полноразмерные датасеты, включая проверку сепарабельности карт Adam,
        обнаруженной в настоящей работе.
  \item Развязанное затухание весов (AdamW)~\cite{loshchilov2019} как отдельное
        сечение: интересно, сохраняется ли сглаживающий эффект при развязке.
  \item Аналитическое описание внутренней границы (между сходимостью и
        ограниченными колебаниями) для двухслойной линейной сети, где
        динамика допускает явный анализ.
\end{enumerate}
